\chapter{Fundamentação}\label{chp:FUNDAMENTACAO}

Este capítulo apresenta os conceitos teóricos necessários para a compreensão da proposta descrita no Capítulo~\ref{chp:PROPOSTA}. Inicialmente, é apresentada a discussão sobre chatbots, abordando sua origem histórica, a definição formal do termo, a taxonomia dos tipos existentes e a justificativa técnica para a escolha de um chatbot baseado em regras neste trabalho. Em seguida, são apresentados os fundamentos de Busca e Recuperação da Informação (BRI), incluindo os principais modelos de recuperação e o papel da indexação e da busca no funcionamento do Elasticsearch, motor utilizado no sistema proposto.

\section{Chatbots}

\subsection{Definição}

A IBM define chatbot como um programa de computador desenvolvido para simular a conversação humana com o usuário final \cite{IBMChatbots}. Essa definição é compatível com a literatura acadêmica sobre o tema. Dale \cite{dale2016chatbots} define chatbot, de forma ampla, como qualquer aplicação de software capaz de manter um diálogo com um ser humano por meio de linguagem natural, independentemente de o sistema ser chamado de assistente digital, interface conversacional ou chatbot. O conceito subjacente é sempre o de alcançar um resultado por meio de uma interação dialógica com uma máquina.

Já Adamopoulou e Moussiades \cite{adamopoulou2020chatbots} caracterizam os chatbots como sistemas de diálogo natural que simulam conversas humanas, podendo ser implementados essencialmente por meio de duas abordagens: a correspondência de padrões (\textit{pattern matching}), em que a resposta é selecionada a partir de regras pré-definidas que casam com a entrada do usuário, e o aprendizado de máquina, em que o sistema aprende a gerar ou selecionar respostas a partir de grandes volumes de dados de treinamento. Essa distinção entre abordagem baseada em regras e abordagem baseada em aprendizado de máquina é a base da taxonomia apresentada a seguir.

\subsection{Origem e Evolução dos Chatbots}\label{sec:origem-evolucao-chatbots}

A história dos chatbots teve início na década de 1960, através do desenvolvimento do programa ELIZA, idealizado por Joseph Weizenbaum no Massachusetts Institute of Technology (MIT). ELIZA foi um chatbot projetado para simular uma conversa com um psicoterapeuta, funcionando com base em regras simples de construção de palavras e scripts pré-definidos \cite{weizenbaum1966eliza}. ELIZA não compreendia de fato o conteúdo das mensagens, mas conseguia simular interações humanas de maneira revolucionária para a época, sendo sua criação um marco do início da exploração de agentes conversacionais no campo da inteligência artificial (IA) e da linguística computacional.

Outro marco na evolução dos chatbots ocorreu em 1972, quando o psiquiatra Kenneth Colby, na Universidade de Stanford, desenvolveu o PARRY, modelo que, diferentemente do ELIZA, simulava o comportamento de um paciente com esquizofrenia paranoide. O PARRY foi submetido a uma versão do Teste de Turing \cite{turing1950}, sendo avaliado por psiquiatras que interagiram com o programa e analisaram suas transcrições. Os especialistas identificaram corretamente cerca de 48\% das trocas, valor que demonstra que, mesmo para profissionais da área, o PARRY se passava, na maioria das situações, por um paciente humano \cite{colby1972parry}. Essa capacidade de simulação evidencia o avanço no uso de agentes conversacionais como ferramentas de investigação psicológica, mesmo em uma época em que a tecnologia não era tão avançada quanto hoje, reforçando o PARRY como uma etapa importante na história da IA conversacional.

Após os primeiros experimentos com o ELIZA e o PARRY, o desenvolvimento de chatbots avançou consideravelmente com a criação do \textit{Artificial Linguistic Internet Computer Entity} (A.L.I.C.E), em 1995, por Richard Wallace. Diferente de seus predecessores, A.L.I.C.E foi projetado com o uso de uma linguagem específica chamada AIML (\textit{Artificial Intelligence Markup Language}), que permitiu a construção de padrões de entrada e respostas de forma modular e escalável, facilitando sua expansão e adaptação para diferentes tipos de diálogos \cite{wallace2009alice}. Mesmo sem utilizar técnicas de aprendizado de máquina, sua estrutura baseada em regras, somada ao grande volume de dados escritos manualmente, possibilitava interações elaboradas, tendo A.L.I.C.E conquistado reconhecimento internacional após vencer três vezes o Prêmio Loebner (2000, 2001 e 2004), competição baseada no Teste de Turing \cite{turing1950} que avalia o grau de humanização de um chatbot \cite{shum2018eliza}.

A partir de 2011, com o lançamento da Siri pela Apple, integrada ao iPhone 4S, iniciou-se uma nova era no desenvolvimento de chatbots inteligentes, marcando a popularização dos assistentes pessoais inteligentes no cotidiano dos usuários. A Siri foi pioneira ao trazer uma interface de voz capaz de compreender comandos naturais e executar tarefas diversas, integrando múltiplas fontes de dados e sensores embarcados no dispositivo \cite{hoy2018}. Ao longo da década seguinte surgiram diversos outros assistentes pessoais inteligentes, como a Cortana da Microsoft e a Alexa da Amazon, ambas em 2014, e o Google Assistant, em 2016. Esses agentes são capazes de integrar informações provenientes de múltiplos sensores, como localização, clima, movimento, toques e gestos, além de acessarem diversas fontes de dados pessoais, como músicas, calendários, e-mails e perfis sociais \cite{shum2018eliza}.

Mais recentemente, a evolução dos chatbots alcançou um novo patamar com a adoção de modelos de linguagem de grande escala, os LLMs (\textit{Large Language Models}) \cite{IBMLLM}. Ao contrário dos chatbots mais tradicionais, que funcionam com base em fluxos fixos ou regras pré-programadas, os LLMs são treinados com enormes volumes de texto por meio de técnicas de \textit{Deep Learning} e conseguem lidar com tarefas bem mais complexas, como responder perguntas abertas, escrever textos longos e bem estruturados, resumir documentos, traduzir idiomas e explicar conceitos técnicos. Essa flexibilidade e capacidade de adaptação a diferentes situações tornam esses modelos especialmente úteis em áreas como educação, atendimento ao cliente, saúde e produtividade no ambiente corporativo \cite{bommasani2021opportunities}. A Seção~\ref{sec:tipos-chatbots} situa essa abordagem dentro da taxonomia de chatbots, e a Seção~\ref{sec:busca-recuperacao-informacao} retoma os LLMs sob a perspectiva da Recuperação da Informação, relacionando-os ao modelo probabilístico apresentado a seguir.

\subsection{Tipos de Chatbots}\label{sec:tipos-chatbots}

A IBM, por meio da publicação \textit{``Seis tipos de chatbots e como escolher o ideal para sua empresa''} \cite{IBMChatbots}, descreve seis categorias de chatbots disponíveis atualmente, apresentadas a seguir.

\subsubsection{Chatbots Baseados em Menus ou Botões}
Representam o nível mais básico entre os chatbots. A interação com o usuário se dá pela seleção de opções em um script que se encaixa melhor nas suas necessidades, podendo haver mais opções conforme as escolhas forem sendo feitas, até chegar à opção mais adequada. Esses chatbots operam como árvores de decisão e são bons para tarefas transacionais, sendo focados na resolução de funcionalidades simples e úteis para responder perguntas repetitivas, previsíveis e diretas. A dificuldade aparece em solicitações menos genéricas, pois as opções de resposta são pré-definidas e limitadas. Caso a necessidade do usuário não esteja disponível no menu, o chatbot se torna inútil.

\subsubsection{Chatbots Baseados em Regras}
Possuem a funcionalidade da árvore de decisão, assim como os chatbots baseados em menu, adicionando também regras de lógica condicional para o desenvolvimento de fluxos de automação de conversa. Atuam essencialmente com FAQs interativas, necessitando de uma quantidade de combinações pré-definidas de perguntas e respostas para a interação com o usuário. Operam com detecção básica de palavras-chave, sendo fáceis de treinar e funcionando bem quando recebem perguntas dentro de um escopo pré-definido, mas têm dificuldade para lidar com perguntas mais complexas e não previstas no escopo.

\subsubsection{Chatbots Impulsionados por IA}
Possuem capacidade de compreensão das perguntas do usuário com recursos de IA e Entendimento de Linguagem Natural (NLU), conseguindo detectar rapidamente as informações relevantes a partir das interações com o usuário, tornando a interação mais conversacional. Possuem a possibilidade de buscar mais informações com o usuário para o entendimento correto do escopo fornecido e solicitado, podendo apresentar opções de resposta entre as quais o usuário escolhe a que mais se encaixa no contexto encontrado. Utilizam algoritmos de aprendizado de máquina em sua implementação, permitindo entendimento, aprendizagem e desenvolvimento contínuo.

\subsubsection{Voice Chatbots}
São ferramentas de conversação que permitem a interação falada como alternativa à interação textual, podendo oferecer as mesmas funcionalidades dos chatbots de IA convencionais, mas implementadas com tecnologia \textit{Text-to-Speech} e \textit{Speech-to-Text}. Com a aplicação do Processamento de Linguagem Natural e interações com tecnologias de computação e telefonia, os chatbots por voz conseguem compreender perguntas faladas, analisar a intenção do usuário e apresentar respostas em tom de conversa, aumentando o engajamento e a imersão.

\subsubsection{Chatbots de IA Generativa}
Oferecem funcionalidades mais avançadas que os anteriores, com maior fluência na compreensão da linguagem comum e maior adaptação à linguagem do usuário, podendo gerar conteúdo novo como resultado sendo texto, imagens ou áudio,com base em \textit{Large Language Models} (LLMs), nos quais são treinados. Apresentam maior personalização aos usuários, aumentando o alcance e a velocidade no atendimento. As interfaces de chatbots com IA generativa podem reconhecer, resumir, traduzir, prever e criar conteúdo em resposta à consulta de um usuário.

\subsubsection{Chatbots Híbridos}
São sistemas de IA conversacional que combinam lógica baseada em regras e recursos baseados em aprendizado de máquina, proporcionando uma experiência mais versátil ao lidar com tarefas de diferentes níveis de complexidade. Funcionam com um conjunto de regras e scripts pré-definidos, que dão estrutura ao diálogo, enquanto a camada de IA contribui com potencial de aprendizado para interações mais complexas.

\subsection{Chatbots Baseados em Regras}

Chatbots baseados em regras operam por meio de um conjunto fechado de padrões de entrada e saída definidos previamente pelo desenvolvedor, sem capacidade de generalização para situações fora desse conjunto. Adamopoulou e Moussiades \cite{adamopoulou2020chatbots} descrevem essa abordagem como \textit{pattern matching}: a entrada do usuário é comparada a padrões pré-estabelecidos e, ao haver correspondência, uma resposta associada a esse padrão é selecionada. Essa lógica está na base de linguagens historicamente usadas para construção de chatbots, como a AIML (\textit{Artificial Intelligence Markup Language}), empregada no A.L.I.C.E \cite{wallace2009alice}, apresentado na Seção~\ref{sec:origem-evolucao-chatbots}.

Singh, Joesph e Jabbar \cite{singh2019rulebased} descrevem um chatbot baseado em regras desenvolvido para responder a consultas administrativas de estudantes de uma universidade (o ``APU Admin Bot''), no qual a correspondência de padrões como palavras, frases e até ações específicas do usuário aciona um conjunto pré-definido de respostas, dispensando a necessidade de programação tradicional ou de grandes volumes de dados de treinamento. Esse caso ilustra a adequação da abordagem baseada em regras a domínios fechados, com escopo de perguntas relativamente previsível, característica também presente no contexto das atas institucionais tratadas neste trabalho.

\subsection{Comparativo entre Tipos de Chatbot e os Requisitos do Projeto}\label{sec:comparativo-chatbots}

A escolha do tipo de chatbot a ser utilizado neste trabalho deve considerar as características do problema apresentado: um sistema de busca em um conjunto fechado de documentos, sendo eles as atas da UFRRJ, sem necessidade de geração de conteúdo novo, com baixo orçamento disponível para licenciamento de ferramentas e que precisa se integrar a um motor de busca textual (Elasticsearch \cite{IBMElasticsearch}) por meio de chamadas a uma API. A Tabela~\ref{tab:comparativo-chatbots} resume, para cada tipo de chatbot apresentado na Seção~\ref{sec:tipos-chatbots}, a adequação a esses requisitos.

\begin{landscape}
\begin{table}[H]
  \centering
  \caption{Comparativo entre tipos de chatbot e os requisitos do projeto}
  \label{tab:comparativo-chatbots}
  \begin{tabular}{p{4cm}p{7.5cm}p{3cm}p{6cm}}
    \hline
    \textbf{Tipo} & \textbf{Adequação ao escopo fechado} & \textbf{Custo de licenciamento} & \textbf{Integração com API externa (Elasticsearch)} \\
    \hline
    Baseado em menus/botões & Limitado: difícil mapear os critérios de busca em menus rígidos & Baixo/nenhum & Possível, mas pouco flexível para parâmetros livres \\
    Baseado em regras & Alto: escopo fechado favorece fluxos com regras condicionais bem definidas & Baixo/nenhum (Botpress v12 open source) & Suporta chamadas HTTP/JS dentro do fluxo \\
    Impulsionado por IA / NLU & Alto poder de generalização, mas exige treinamento de modelos de NLU e maior esforço de configuração & Médio/alto (modelos proprietários) & Possível, porém com custo adicional \\
    Voice chatbot & Não se aplica ao caso de uso (interação textual) & Médio/alto & Possível, mas fora do escopo \\
    IA generativa & Alta flexibilidade, porém risco de respostas não fundamentadas nas atas (alucinação) e custo recorrente de uso de LLM & Alto (APIs de LLM pagas) & Possível, mas exige camada adicional de RAG \\
    Híbrido & Combina regras e IA, mas adiciona complexidade de manutenção desnecessária para o escopo atual & Médio & Possível \\
    \hline
  \end{tabular}
\end{table}
\end{landscape}

Diante desses requisitos, optou-se pelo chatbot baseado em regras, implementado na plataforma Botpress v12 \cite{botpress-about, botpress-docs}. Essa escolha decorre de três fatores combinados: (i) o escopo do problema é fechado e o usuário busca atas dentro de um conjunto de critérios pré-definidos (departamento, disciplina, docentes, discente, data, processo/edital/ata, progressão funcional e outros assuntos), o que se encaixa em um fluxo de diálogo guiado por regras; (ii) o Botpress v12 é uma ferramenta de código aberto, sem custo de licenciamento, podendo ser hospedada localmente \cite{botpress-about}; e (iii) a plataforma permite a execução de ações em JavaScript dentro do fluxo de conversa, viabilizando a integração com a API REST do indexador Node.js, que por sua vez consulta o Elasticsearch \cite{IBMElasticsearch}. Chatbots impulsionados por IA, de IA generativa ou híbridos poderiam oferecer maior flexibilidade na interpretação de perguntas livres, mas introduziriam custo de licenciamento de modelos de linguagem e risco de geração de respostas não embasadas no conteúdo das atas.

\section{Busca e Recuperação da Informação}\label{sec:busca-recuperacao-informacao}

\subsection{Definição}

Recuperação de Informação (RI), do inglês \textit{Information Retrieval}, é definida por Manning, Raghavan e Schütze \cite{manning2008ir} como a atividade de localizar, dentro de uma grande coleção, material de natureza não estruturada que satisfaça uma necessidade de informação expressa por uma consulta. Baeza-Yates e Ribeiro-Neto \cite{baezayates2010modern} complementam essa definição ao descrever a RI como a área responsável por representar, armazenar, organizar e acessar itens de informação de forma a facilitar o acesso do usuário às informações de seu interesse, lidando com a representação, o armazenamento e a recuperação de informação textual.

Diferentemente de uma consulta a um banco de dados estruturado, em que a correspondência entre consulta e registros é exata, a RI trata do problema da relevância: dado um conjunto de documentos e uma consulta, o sistema deve estimar quais documentos são relevantes para a necessidade de informação do usuário, ordenando os resultados por algum critério de pontuação \cite{manning2008ir}.

\subsection{Modelos de Recuperação da Informação}

Baeza-Yates e Ribeiro-Neto \cite{baezayates2010modern} organizam os modelos clássicos de RI em três grandes famílias: o modelo booleano, o modelo vetorial (espaço vetorial) e o modelo probabilístico.

\subsubsection{Modelo Booleano}
Trata-se de um modelo de recuperação de informação baseado na teoria dos conjuntos e na álgebra booleana, neste modelo documentos e consultas são representados como conjuntos de termos combinados por operadores lógicos (\textit{AND}, \textit{OR}, \textit{NOT}). Um documento é considerado relevante se satisfaz exatamente a expressão booleana da consulta, sem que exista uma noção de grau de relevância, ou seja, o documento ou corresponde ou não corresponde a busca realizada.\cite{baezayates2010modern}.

\subsubsection{Modelo Vetorial (Espaço Vetorial)}\label{sec:modelo-vetorial}
No modelo de espaço vetorial, proposto originalmente por Salton, Wong e Yang \cite{salton1975vector}, documentos e consultas são representados como vetores em um espaço de $n$ dimensões, em que cada dimensão corresponde a um termo do vocabulário. O grau de similaridade entre uma consulta e um documento é calculado por uma medida de distância ou similaridade entre os vetores, permitindo ordenar os documentos por um grau contínuo de relevância, em vez da resposta binária do modelo booleano. Esse modelo é a base conceitual tanto de esquemas de ponderação clássicos, como o TF-IDF, quanto das abordagens modernas de busca semântica baseadas em \textit{embeddings}, nas quais documentos e consultas são representados como vetores densos aprendidos por modelos de linguagem, e a similaridade vetorial captura relações semânticas entre termos que vão além da correspondência exata de palavras.

\subsubsection{Modelo Probabilístico}
O modelo probabilístico estima a probabilidade de um documento ser relevante para uma consulta, ordenando os resultados de acordo com essa estimativa \cite{baezayates2010modern}. A família de funções de ranqueamento BM25 (\textit{Best Matching 25}) é a realização mais difundida desse modelo, derivada do arcabouço probabilístico de relevância desenvolvido por Robertson e colaboradores \cite{robertson2009bm25}. O BM25 pondera a frequência de cada termo da consulta no documento, a frequência inversa do termo na coleção (termos raros contribuem mais para a pontuação que termos comuns) e o comprimento do documento, produzindo uma pontuação de relevância para cada documento em relação à consulta. É esse o algoritmo utilizado pelo Elasticsearch para o ranqueamento padrão de resultados \cite{IBMElasticsearch}.

A noção de estimativa de probabilidade que sustenta o modelo probabilístico de RI também é o princípio por trás dos LLMs (\textit{Large Language Models}) que impulsionam os chatbots de IA generativa apresentados na Seção~\ref{sec:tipos-chatbots}. Segundo a IBM \cite{IBMLLM}, um LLM é um modelo de aprendizado de máquina aplicado à IA generativa, treinado com grandes volumes de texto para compreender e gerar linguagem de forma contextualmente coerente. Do ponto de vista matemático, um LLM estima, a cada passo, a probabilidade de cada token do vocabulário ser o próximo em uma sequência, dado o contexto precedente; a representação vetorial desse contexto (\textit{embeddings}) segue a mesma lógica de espaço vetorial discutida na Seção~\ref{sec:modelo-vetorial}. Há, portanto, uma proximidade matemática entre os dois modelos, pois ambos se apoiam em estimativas de probabilidade, ainda que aplicados a fins distintos: o modelo probabilístico de RI estima a relevância de um documento para uma consulta, com o objetivo de ranqueamento, enquanto o LLM estima a probabilidade do próximo token em uma sequência, com o objetivo de geração de linguagem.

\subsubsection{Situando o BM25 e a Busca Semântica}
No sistema proposto neste trabalho, a recuperação textual sobre o conteúdo das atas é realizada pelo Elasticsearch, que utiliza o BM25 como função de ranqueamento padrão \cite{IBMElasticsearch}. A busca por similaridade semântica, já mencionada neste capítulo como mecanismo complementar ao BM25, situa-se dentro da família de modelos vetoriais: tanto a consulta quanto os documentos podem ser representados como vetores de \textit{embeddings}, e a relevância é estimada pela proximidade entre esses vetores no espaço semântico, capturando relações de significado que o BM25, por operar sobre a correspondência lexical de termos, não captura diretamente.

\subsection{O Problema da Busca: Indexação e Recuperação}\label{sec:problema-busca}

O funcionamento de um sistema de RI pode ser dividido em duas etapas principais: a indexação e a busca propriamente dita \cite{manning2008ir}.

A indexação consiste em processar previamente a coleção de documentos para construir uma estrutura de dados que viabilize buscas eficientes sem a necessidade de varrer o conteúdo integral de cada documento a cada consulta. No caso do Elasticsearch, essa estrutura é o índice invertido, no qual cada termo do vocabulário é associado à lista de documentos em que ocorre, junto a estatísticas usadas pelo BM25 (frequência do termo, frequência do documento, comprimento do documento) \cite{IBMElasticsearch}. No sistema proposto, essa etapa corresponde ao processo de extração de texto dos PDFs das atas e ao envio desse conteúdo ao Elasticsearch para indexação, associando a cada documento campos estruturados como título e data da ata, além do conteúdo textual completo.

A busca, por sua vez, consiste em receber uma consulta, processá-la e utilizá-la para percorrer o índice em busca dos documentos mais relevantes, retornando-os ordenados por pontuação. No sistema proposto, a consulta enviada ao Elasticsearch é montada a partir dos critérios informados pelo usuário através do chatbot, e os documentos retornados são ordenados de forma decrescente pelo \textit{score} atribuído pelo BM25, sendo selecionados os três melhores resultados para apresentação ao usuário. O Capítulo~\ref{chp:PROPOSTA} discute, em nível conceitual, como o chatbot resolve a desconexão entre o vocabulário livre do usuário e a organização do índice na geração dessa consulta; o Capítulo~\ref{chp:EXPERIMENTOS} apresenta os detalhes de implementação dessa etapa.
